% Standalone export: рис. 2.11 — компонентная диаграмма архитектуры (для презентации)
\documentclass[tikz,border=6pt]{standalone}
\usepackage[utf8]{inputenc}
\usepackage[T2A]{fontenc}
\usepackage[russian]{babel}
\usetikzlibrary{shapes.geometric, arrows.meta, positioning}

\begin{document}
\begin{tikzpicture}[
    scale=0.95, transform shape,
    every node/.style={font=\small},
    layer/.style={draw, thick, rounded corners=6pt,
      minimum width=3.4cm, minimum height=1.2cm,
      align=center},
    store/.style={draw, thick, cylinder, shape border rotate=90,
      minimum width=2.8cm, minimum height=1.2cm, aspect=0.2,
      align=center},
    arrow/.style={thick,->,>=stealth},
    label/.style={font=\footnotesize, midway}
]

\node (react) [layer, fill=white] at (0, 0)
  {\textbf{Веб-клиент (SPA)}};

\node (api) [layer, fill=white] at (6, 0)
  {\textbf{REST API сервер}};

\node (pg) [store, fill=white] at (12.5, 1.5)
  {\textbf{Реляционная СУБД}};

\node (redis) [store, fill=white] at (12.5, -1.5)
  {\textbf{In-memory хранилище}};

\draw [arrow] ([yshift=5pt]react.east) -- node[label, above] {HTTP / REST} ([yshift=5pt]api.west);
\draw [arrow, dashed] ([yshift=-5pt]api.west) -- node[label, below] {JSON} ([yshift=-5pt]react.east);

\draw [arrow] ([yshift=8pt]api.east) -- ++(0.8,0) |- ([yshift=5pt]pg.west)
  node[pos=0.8, above, font=\scriptsize] {запрос к БД};
\draw [arrow, dashed] ([yshift=-5pt]pg.west) -| ++(-0.6,0) |- ([yshift=2pt]api.east);

\draw [arrow] ([yshift=-2pt]api.east) -- ++(0.8,0) |- ([yshift=5pt]redis.west)
  node[pos=0.8, above, font=\scriptsize] {ключ--значение};
\draw [arrow, dashed] ([yshift=-5pt]redis.west) -| ++(-0.6,0) |- ([yshift=-8pt]api.east);

\node [font=\footnotesize\itshape, text=gray] at (0, -1.8) {Уровень представления};
\node [font=\footnotesize\itshape, text=gray] at (6, -1.8) {Уровень приложения};
\node [font=\footnotesize\itshape, text=gray] at (12.5, -2.8) {Уровень данных};
\end{tikzpicture}
\end{document}
